\subsection{Análisis del Módulo 3}
hola ff ff
\subsubsection{Modelo de base de datos}
% Tabla comparativa entre SQL y NoSQL
\begingroup
\setlength{\tabcolsep}{6pt} % Espacio entre columnas
\renewcommand{\arraystretch}{1.5} % Altura de las filas

\begin{longtable}{|p{3.5cm}|p{5.5cm}|p{5.5cm}|}
\caption{Modelo de base de datos.}
\label{tab:modelodb} \\
\hline
\textbf{Característica} & \textbf{SQL (Relacional)}\cite{Oppel2011} & \textbf{NoSQL (No Relacional)}\cite{Gomez} \\
\hline
\textbf{Modelo de Datos} & \textbf{Esquema Fijo} (Tablas, Filas y Columnas predefinidas en una estructura). & \textbf{Esquema Dinámico/Flexible} (Documentos, Pares Clave-Valor, Grafos, etc.). \\
\hline
\textbf{Escalabilidad} & Escalabilidad  principalmente Vertical (mejorar servidor) y horizontal (sharding). & Principalmente permite una escalabilidad horizontal (añadir más servidores o nodos), lo que permite mayor distribución \\
\hline
\textbf{Transacciones } & ACID (Atomicidad, Consistencia, Aislamiento, Durabilidad). & BASE (Básicamente Disponible, Estado Suave, Consistencia Eventual). \\
\hline
\textbf{Arquitectura} & Monolítica, centralizada en un único servidor o clúster. & Distribuida, diseñada para clústeres de nodos a lo largo de múltiples servidores. \\
\hline
\textbf{Ventajas} & Basado en el modelo relacional permite la integridad de datos,garantizada, consultas complejas con JOINs eficientes y modelado de relaciones claro. & Permite una alta velocidad y consistencia en inserción/lectura, adaptabilidad a cambios, y rendimiento en grandes volúmenes de datos. \\
\hline
\textbf{Desventajas} & Bastante robustas lo que las vuelve menos flexibles a cambios en el esquema. Teniendo un rendimiento limitado para escrituras y lecturas distribuidas. & Compromete la fiabilidad a cambio de rendimiento, y al tener no tener propiedades ACID tan estrucuturadas no permite relaciones complejas. \\
\hline
\end{longtable}
\endgroup

\subsubsection{Base de datos NoSQL}
\begin{landscape}
\begingroup
\setlength{\tabcolsep}{.1pt} % Espacio entre columnas reducido
\renewcommand{\arraystretch}{1.5} % Altura de las filas ligeramente reducida

\begin{table}[h!]

\caption{Bases de datos NoSQL}
\label{tab:nosql}
\footnotesize 
\begin{tabular}{|>{\raggedright\arraybackslash}p{2.5cm}|>{\raggedright\arraybackslash}p{2.5cm}|>{\raggedright\arraybackslash}m{2.5cm}|>{\raggedright\arraybackslash}p{4cm}|>{\raggedright\arraybackslash}p{2.5cm}|>{\raggedright\arraybackslash}p{2.5cm}|>{\raggedright\arraybackslash}p{2.5cm}|>{\raggedright\arraybackslash}p{3cm}|}
\hline
\textbf{Base} & \textbf{Modelo de Datos} & \textbf{CAP}& \textbf{Características} & \textbf{Licencia} & \textbf{Propiedades ACID} & \textbf{Propiedades BASE} & \textbf{Implementación en la Nube o Local} \\
\hline
\textbf{MongoDB}\cite{zotero-item-107} & Documento (nativo BSON/JSON) & CP & Permite esquemas flexibles, asignando cada documento a un objeto. & Código abierto SSPL/Apache & Media:Permite transacciones ACID entre documentos pero de forma limitada. & Fuerte (Consistencia Eventual en clústeres distribuidos). & Nube (MongoDB Atlas, AWS, GCP, Azure) y Local\\
\hline
\textbf{Cassandra}\cite{zotero-item-108} & Columna Amplia (Wide-Column) & AP & Permite una escalabilidad lineal al optimiza las consultas por medio de CQL(Lenguaje de consulta de Cassandra) para grandes cantidades de datos. & Apache & Bajo (Garantías mínimas de consistencia y aislamiento). & Fuerte (Alta disponibilidad y particionamiento automático).& Nube (AWS, GCP, Azure) y Local \\
\hline
\textbf{Redis}\cite{zotero-item-111} & Clave-Valor (Key-Value) & CP & Almacenamiento en memoria (*in-memory*), uso primario como caché y broker de mensajes. & BSD & Parcial (Transacciones atómicas básicas en comandos únicos). & Fuerte (Alta velocidad y disponibilidad por naturaleza *in-memory*). & Nube (ElastiCache, Memorystore) y Local\\
\hline
\textbf{InfluxDB}\cite{zotero-item-113} & Series de Tiempo (TSDB) & AP & Optimizado para ingesta masiva de datos indexados por secuencias de tiempo. & MIT y Comercial & Bajo (Prioriza la velocidad de ingesta sobre la atomicidad). & Fuerte (Alta disponibilidad, optimizado para particionamiento temporal). & Nube (InfluxDB Cloud) y Local\\
\hline
    \end{tabular}
\end{table}
\endgroup
\end{landscape}